\section{Introduction}
\label{sec:introduction}

Client endpoints represent the most vulnerable perimeter in modern web and mobile application security~\cite{syed2022zero}. While identity providers (IdPs) enforce robust server-side cryptography, user interactions terminate on client devices that frequently operate outside corporate control~\cite{jannett2024sok}. Bring Your Own Device (BYOD) policies, shared public workstations, and commodity operating systems expose users to host-level threats, including credential-harvesting malware, memory scrapers, and malicious browser extensions~\cite{gautam2025passwords, lee2024beneath}. In this threat environment, traditional authentication factors fail. Passwords and one-time codes transmitted across compromised input devices are recorded instantly by low-level keyloggers~\cite{amft2023weve}.

To counter phishing and network relay, industry standards have converged on hardware-anchored credentials, specifically FIDO2/WebAuthn~\cite{angelogianni2024how, nist80063b4}. WebAuthn deterministically halts network-level Adversary-in-the-Middle (AitM) proxies through public-key cryptography bound to the browser origin~\cite{golaszewski2026binding}. However, WebAuthn assumes an uncompromised client operating environment. When malware or reverse-proxy toolkits target the session lifecycle, adversaries intercept bearer session tokens post-authentication or exploit UI overlay deception on client devices~\cite{mahdad2024breaching, hilbig2023protect, lim2025phishers}. Furthermore, hardware security keys face physical distribution, recovery, and user adoption barriers in consumer web applications and public terminal settings~\cite{kepkowski2023challenges}.

This paper proposes the \textit{Blind Cognitive Multi-Factor Authentication} (BC-MFA) architecture. BC-MFA addresses step-up authentication on untrusted web and mobile endpoints without requiring external cryptographic hardware tokens. The core mechanism separates the visual stimuli displayed on the untrusted screen from the actual secret verified by the IdP. The server issues a dynamic cognitive challenge: an ephemeral, quadrant-balanced visual transformation matrix. The user mentally combines the challenge with a pre-shared cognitive secret (a private anchor quadrant and a non-zero toroidal geometric offset path) to generate a one-time response~\cite{binbeshr2025challenge}. An adversary capturing keystrokes or recording the display learns neither the underlying secret nor the rule governing subsequent challenges. To mitigate automated bot replay and remote human solver delegation, BC-MFA couples this cognitive ceremony with passive behavioral anomaly telemetry capturing cognitive hesitation latency, pointer or touchscreen kinematics, and keystroke cadence~\cite{incel2021dakota, khan2024mouse}.

Three concrete research questions guide this inquiry. First, \textbf{RQ1} examines whether quadrant-balanced challenge-response permutations over toroidal grids can guarantee bounded combinatorial ambiguity against continuous display capture across $k$ compromised sessions while keeping expected human mental latency within practical bounds ($\le 5\text{ s}$). Second, \textbf{RQ2} investigates the protocol bindings required to integrate cognitive verification, RFC 7636 PKCE, and RFC 9449 Demonstrating Proof-of-Possession (DPoP) access token binding into standard OpenID Connect Core 1.0 public-client flows using registered RFC 8176 Authentication Method Reference (\texttt{amr}) claims without requiring mandatory relying party modifications. Third, \textbf{RQ3} evaluates whether passive interaction kinematics (cognitive hesitation, pointer/touch trajectories, and short-input keystroke cadence) provide sufficient discriminatory signal to separate authentic human operators from automated replay bots and off-device relay solvers on heterogeneous endpoints.

This paper delivers four core contributions. First, it defines the formal protocol specification and worked toroidal transformation mechanics for Blind Cognitive MFA on untrusted web and mobile public clients. Second, it specifies a standards-compliant federation profile mapping cognitive, behavioral, PKCE, and DPoP key-binding ceremonies into OpenID Connect Core 1.0 using RFC 7636, RFC 8176, and RFC 9449. Third, it provides a mathematical threat evaluation (covering keyloggers, screen scrapers, SIM swapping, and AitM proxies, while demarcating first-contact active reverse-proxy limits) validated by a reproducible 500-trial rational-adversary Monte Carlo intersection simulation. Fourth, it establishes a literature-grounded latency model and a formal, ethics-compliant empirical evaluation plan for the Master's dissertation.

Section~\ref{sec:methodology_and_prisma} details the systematic literature review methodology. Section~\ref{sec:related_work} positions this proposal relative to existing cognitive and biometric literature. Section~\ref{sec:architecture} details the system architecture, worked toroidal matrix transformations, and OpenID Connect token bindings. Section~\ref{sec:security_analysis} provides formal threat modeling and Monte Carlo intersection analysis. Section~\ref{sec:usability_and_overhead} analyzes cognitive latency bounds, fault tolerance, and accessibility. Section~\ref{sec:experimental_design_and_thesis} formalizes the empirical dissertation evaluation framework, and Section~\ref{sec:conclusion} concludes the paper.
